%webamc CODE=sae12/2025/routage
%webamc TITLE=Routage IP

On considère dans cet exercice le réseau suivant. Les numéros entre
crochets correspondent à des numéros d'interface (p.ex., $[0]$ pour
eth0).\\
\begin{center}
  \begin{minipage}{.67\textwidth}
    \scalebox{0.92}{\begin{tikzpicture}[node distance=2.3cm]
  \small
  \routeur{}{R1}{90}{R1}
  \lan{left of=R1}{L2}{réseau $1.2.0.0/16$}
  \lan{right of=R1}{L3}{réseau $1.3.0.0/16$}
  \routeur{right of=L3}{R2}{90}{R2}
  \lan{below of=R1,node distance=1.8cm}{L1}{réseau $1.1.0.0/16$}
  \machine{left of=L1}{PC1}{90}{PC1}
  \lan{below of=R2,node distance=1.8cm}{L4}{réseau $1.4.0.0/16$}
  \internet{right of=R2}{I}
  %
  \draw[cable,-] (PC1) -- node[pos=0.2,above] {\scriptsize $[0]$} (L1);
  \draw[cable,-] (R1) -- node[pos=0.2,left] {\scriptsize $[1]$} (L1);
  \draw[cable,-] (R1) -- node[pos=0.2,above] {\scriptsize $[0]$} (L2);
  \draw[cable,-] (R1) -- node[pos=0.2,above] {\scriptsize $[2]$} (L3);
  \draw[cable,-] (R2) -- node[pos=0.2,above] {\scriptsize $[1]$} (L3);
  \draw[cable,-] (R2) -- node[pos=0.2,left] {\scriptsize $[2]$} (L4);
  \draw[cable,-] (R2) -- node[pos=0.2,above] {\scriptsize $[0]$} (I);
\end{tikzpicture}
}
  \end{minipage}
  \begin{minipage}{.31\textwidth}
    \begin{center}
      \begin{tabular}{|c|c|c|}
        \hline
        Hôte & Interface & Adr. IP\\
        \hline
        \hline
        PC1 & eth0 & 1.1.0.1\\
        \hline
        R1 & eth0 & 1.2.255.254\\
        & eth1 & 1.1.255.254\\
        & eth2 & 1.3.255.254\\
        \hline
      \end{tabular}
    \end{center}
  \end{minipage}
\end{center}

\noindent%
  \begin{minipage}{.49\textwidth}
On donne ci-contre la table de routage de PC1. On
suppose que les tables de routage permettent à PC1 de communiquer avec
tous les réseaux (internet compris).
  \end{minipage}
  \begin{minipage}{.02\textwidth}\end{minipage}
  \begin{minipage}{.46\textwidth}

  \begin{tabular}{|c|c|c|c|}
    \hline
    Destination & Masque & Passerelle & Interface\\
    \hline
    \hline
    {\em C1} & 255.255.0.0 & {\em C2} & eth0\\
    0.0.0.0 & 0.0.0.0 & {\em C3} & {\em C4}\\
    \hline
  \end{tabular}
  \end{minipage}
\\
